\section{Reconstruction of positive geometry from nonadic memory}
\label{sec:memory}

Dodecaphasic extraction preserves more information than a deformation
direction. The direction selects a motion; the full record
also determines the partition on which that motion acts. Recovering
the record from its memories allows this distinction to be
formulated as a reconstruction theorem. The result follows from the
composition of the nonadic compressions of the corpus \cite{hmtX}; its geometric
application preserves the uniform charge and the complements of both stages.

In this section, coordinates are numbered from zero to eleven. Let
\(\mathsf S:\RR^{12}\to\RR^{12}\) be the shift
\((\mathsf Sx)_i=x_{i+1}\), with indices modulo twelve. In the Euclidean chart
subsequent to the enriched state, define
\[
 T_j=\frac{8I+\mathsf S^j}{9},\qquad
 D_j=\frac{\sqrt8}{9}(\mathsf S^j-I),\qquad j\in\{3,4\}.
\]
The coefficient eight counts the ordinary positions of the nonadic chart;
the ninth position contains transport. Orthogonality of the shift
implies, by direct multiplication,
\[
 T_j^*T_j+D_j^*D_j=I.
\]
The first memory is the complement of the first compression; the second
is taken on the state that has already passed through that compression. Thus,
\begin{equation}
 q=\mathbf1^*K,\qquad m_0=D_3K,\qquad m_1=D_4T_3K,
 \qquad MK=(m_0,m_1).
 \label{eq:mem-data}
\end{equation}
The order of composition is part of the object. The map
\(M:\RR^{12}\to\RR^{24}\) preserves two correlated observations,
produced by the same linear trajectory.

\begin{theorem}[Exact inversion of the memory readout]
\label{thm:memory-inversion}
The map \(K\mapsto(q,MK)\) is a linear isomorphism between
\(\RR^{12}\) and \(\RR\times\operatorname{im}M\). Its inverse is obtained
through
\begin{align}
 T_3^{-1}&=\frac{512I-64\mathsf S^3+8\mathsf S^6-\mathsf S^9}{455},
 \label{eq:memory-T-inverse}\\
 b&=-\frac9{\sqrt8}m_0,\qquad
 c=-\frac9{\sqrt8}T_3^{-1}m_1,\qquad w_i=c_i-b_{i+1},\notag\\
 B_0&=0,\qquad B_i=\sum_{j=0}^{i-1}w_j,\qquad
 K_i=\frac{q+\sum_{j=0}^{11}B_j}{12}-B_i.
 \label{eq:memory-inverse}
\end{align}
\end{theorem}
\begin{proof}
Writing \(X=\mathsf S^3\), one has \(X^4=I\) and
\((8I+X)(512I-64X+8X^2-X^3)=4095I\). Dividing by nine
proves \eqref{eq:memory-T-inverse}. All operators used are
polynomials in the same shift and commute. Hence,
\[
 b=(I-\mathsf S^3)K,\qquad c=(I-\mathsf S^4)K,
 \qquad w_i=K_i-K_{i+1}.
\]
Telescoping gives \(B_i=K_0-K_i\); summing this equality over
the twelve positions yields \eqref{eq:memory-inverse}. Furthermore,
\(MK=0\) is equivalent to invariance under \(\mathsf S^3\) and
\(\mathsf S^4\). Since \(\mathsf S=\mathsf S^4(\mathsf S^3)^{-1}\),
the kernel is exactly \(\RR\mathbf1\). The charge recovers this
component and proves injectivity. The decomposition
\(K=(q/12)\mathbf1+K^\perp\) proves surjectivity onto the
declared codomain.
\end{proof}

The geometric form is determined by information that a compressed
reading appeared to have displaced. The vector complement contains the
differences between positions; the charge contains the common component. Their
combination reconstructs the twelve coordinates before normalizing them as
lengths. This assertion relates effective operations: compression,
recording of complements and inversion. Conservation of a norm expresses
a distinct property, which by itself is weaker than this recoverability.

\subsection{Compatible memory cone and stability of minors}

Let \(L\) be the pseudoinverse of \(M\), with image
\(\mathbf1^\perp\). Equivalently,
\[
 L=\left(M^*M\big|_{\mathbf1^\perp}\right)^{-1}M^*,
 \qquad LM=P_{11},\qquad K=\frac q{12}\mathbf1+Lm.
\]
The restricted inverse notation is interpreted as an operator taking
values in \(\mathbf1^\perp\). Strict positivity of the record
is equivalent to membership in the open cone
\begin{equation}
 \mathcal M_+=\left\{(q,m)\in\RR\times\operatorname{im}M:
 \frac q{12}\mathbf1+Lm>0\right\}.
 \label{eq:memory-positive-cone}
\end{equation}
The inequalities are coordinatewise. The cone has dimension twelve; its section
at a fixed positive charge has dimension eleven. In particular, the twenty-four
entries of the two memories satisfy compatibility relations.

\begin{proposition}[Optimal reconstruction constant in Euclidean norm]
\label{prop:memory-stability}
One has \(\|L\|=9/4\). For arbitrary perturbations of the observations,
\begin{equation}
 \|\delta K\|^2\leq\frac{|\delta q|^2}{12}
 +\frac{81}{16}\bigl(\|\delta m_0\|^2+\|\delta m_1\|^2\bigr).
 \label{eq:memory-stability}
\end{equation}
\end{proposition}
\begin{proof}
The Fourier basis diagonalizes \(\mathsf S\). If \(z^{12}=1\),
the eigenvalue of \(M^*M\) in the corresponding mode is
\[
 \lambda(z)=\frac8{81}|z^3-1|^2
 +\frac8{81}|z^4-1|^2\left|\frac{8+z^3}{9}\right|^2.
\]
Evaluating all twelve modes gives
\[
 \operatorname{Spec}(M^*M)=
 \left\{0^{[1]},(16/81)^{[2]},(8/27)^{[2]},(32/81)^{[1]},
 (952/2187)^{[4]},(1256/2187)^{[2]}\right\}.
\]
Bracketed exponents indicate multiplicity. The smallest positive
eigenvalue is \(16/81\), giving \(\|L\|=9/4\). The uniform component
\((\delta q/12)\mathbf1\) is orthogonal to \(L\delta m\); the Pythagorean
identity and the norm of \(L\) yield \eqref{eq:memory-stability}.
\end{proof}

To pass to interval geometry, return to numbering from one to twelve
and set \(d_j=K_j/q\), \(t_0=0\),
\(t_j=\sum_{r=1}^j d_r\). The cone \eqref{eq:memory-positive-cone}
determines exactly \(0=t_0<t_1<\cdots<t_{12}=1\), and therefore
\[
 \Delta_{ij}=t_j-t_i=\sum_{i<r\leq j}d_r>0\qquad(i<j).
\]
These are the minors of the path matrix constructed in Section
\ref{app:network-cluster}. Their dependence on memory is
explicit. The points \((t_j,t_j^2)\) are likewise determined and
produce the rank-one amplituhedral realization developed below.

If the square root of the right-hand side of \eqref{eq:memory-stability} is smaller than
\(\min_j K_j\), the perturbed reconstruction remains in the positive cone.
For the resulting record, \(\min_jK_j=58\); at fixed charge it suffices that
\[
 \sqrt{\|\delta m_0\|^2+\|\delta m_1\|^2}<\frac{232}{9}.
\]
The bound preserves vertex order and the sign of every ordered
minor. It describes stability of this mathematical chart, with the
stated normalization. For observations outside \(\operatorname{im}M\),
reconstruction uses the projection \(MLm\); the residue
\((I-ML)m\) is retained separately when recovery of the full
observation is required.

\subsection{Nodal recovery and boundary reduction}

The resulting partition also has a quadratic realization:
\[
 \mathcal E_d(f)=\sum_{j=1}^{12}\frac{|f_j-f_{j-1}|^2}{d_j}
 =f^*L_df.
\]
The nodal Laplacian recovers each length through
\(d_j=-1/(L_d)_{j-1,j}\). If only the global
endpoints are retained, minimization eliminates the interior nodes and instead
produces the matrix \(\left(\begin{smallmatrix}1&-1\\-1&1\end{smallmatrix}\right)\),
since \(\sum_jd_j=1\). The nodal reader distinguishes the interval form;
the two-endpoint reader observes only its total length.

The proof of this elimination is given in full in the energy
section: the minimizing affine prolongation \(Hf\) and the detail \(z\),
which vanishes at the old nodes, satisfy
\[
 \mathcal E_{\widetilde d}(Hf+z)
 =\mathcal E_d(f)+\mathcal E_{\widetilde d}(z).
\]
The second term is nonnegative and vanishes exactly for zero
detail. The effective reading preserves the minimum energy; the pair \((f,z)\)
also preserves the refined configuration. Thus, positivity of minors,
energy positivity and memory recoverability are related through
maps defined on the same partition, retaining their respective
target spaces.
